
\begin{remark}
The scaling in the proof of Theorem~\ref{thm:main} turns Proposition~\ref{prop:halfpotential} into $\sup_{a\in\R}\E\abs{Q-a}^{-1/2}\le77\,K/\sqrt s$ whenever $s_3(A)\ge s>0$. One large eigenvalue does not suffice for such a bound: if $X_1$ is uniform on $[0,1]$, then $\E\abs{X_1^2}^{-1/2}=\E X_1^{-1}=\infty$.
\end{remark}

\begin{remark}
The number five is used only in the case without pivots, through Proposition~\ref{prop:signs} with $m=3$; the diagonal pivot uses four large eigenvalues of $A$ and the pair cases use three. Four eigenvalues without pivots need not give a third large singular value in any off-diagonal block. Let $N>256$, let $w\in\R^N$ have all entries equal to $N^{-1/2}$, and let $A$ be the direct sum of two copies of the $(N+1)\times(N+1)$ matrix $\begin{pmatrix}0&w^T\\w&0\end{pmatrix}$. Its nonzero eigenvalues are $1,1,-1,-1$, its diagonal vanishes and every principal $2\times2$ minor has modulus at most $1/N<\zeta$, so $A$ has no pivot. For every partition $[n]=I\sqcup J$ the part of each copy in $A_{I,J}$ lies in a single row or a single column, namely that of the centre of the copy, so $A_{I,J}$ has rank at most two and $E_3(A_{I,J})=0$.
\end{remark}

